\documentclass[11pt]{article}
\usepackage[T1]{fontenc}
\usepackage{lmodern}
\usepackage[margin=1in]{geometry}
\usepackage{amsmath,amssymb,amsthm}
\usepackage{booktabs}
\usepackage[round,authoryear]{natbib}
\usepackage{microtype}
\usepackage[colorlinks=true,allcolors=blue]{hyperref}

\newtheorem{theorem}{Theorem}[section]
\newtheorem{proposition}[theorem]{Proposition}
\newtheorem{lemma}[theorem]{Lemma}
\theoremstyle{remark}
\newtheorem{remark}[theorem]{Remark}
\DeclareMathOperator{\diag}{diag}
\DeclareMathOperator{\Cov}{Cov}
\DeclareMathOperator{\Var}{Var}
\DeclareMathOperator{\sign}{sign}
\newcommand{\E}{\mathbb{E}}
\newcommand{\Prob}{\mathbb{P}}
\newcommand{\R}{\mathbb{R}}
\newcommand{\soft}{S}

\title{
  Learning Rare Features: Normalization, Stochastic Optimization, and Early
  Stopping
}
\author{Johan Larsson}
\date{Working manuscript --- September 16, 2026}

\begin{document}
\maketitle

\begin{abstract}
  This working investigation asks how normalization changes the learning of rare
  binary features under a finite training budget. For a fixed block of common
  predictors and one rare binary feature, we derive a population gradient-flow
  limit in which incomplete common-feature learning contributes to the rare
  coefficient even as common--rare Pearson correlations vanish. Exact
  finite-update formulas for gradient descent and fresh-sample stochastic
  gradient descent separate initial uncertainty, multiplicative update variance,
  and observation noise. Under a bounded common-block assumption, an exact
  second-moment operator criterion yields stability thresholds across small,
  intermediate, and large products of batch size and rare prevalence. The
  results use known population centering and fixed scales. For two binary
  features and prevalence $q_n = n^{-1/2}$, a specified decreasing-step schedule
  using $n$ fresh observations attains rare-coefficient mean squared error
  asymptotic to $4\sigma^2/(3nq_n)$. A controlled fixed-data pilot finds
  substantial rare-subgroup prediction differences whose direction changes with
  signal, dependence, batch size, and budget under separate oracle step-size
  tuning. A separate sampled-data and sparse categorical study finds that
  scaling effects remain under estimated preprocessing and validation selection,
  while missing rare validation observations can conceal explosive subgroup
  error. Broader finite-sample guarantees, adaptive-optimizer comparisons, and
  generalization beyond the categorical benchmark remain research objectives.
  Earlier weighted-lasso and ridge results are retained as
  explicit-regularization baselines; no standalone novelty claim is made for
  these calculations.
\end{abstract}

\section{Motivation and scope}

Feature scaling sets relative learning rates as well as relative penalties.
For a rare binary feature, scaling by its standard deviation removes the
small marginal variance from the population Hessian, but it also increases
the magnitude of individual observations. Dependence with common features
can further change the learning trajectory. We ask when normalization helps
learn rare features under a finite training budget, and when stochastic
updates or dependence offset that benefit.

The connection between early stopping and regularization is established:
\citet{ali2019early} compare the risks of least-squares gradient flow and
ridge regression. Optimization geometry also determines implicit bias in
underdetermined models \citep{gunasekar2018geometry}, and preconditioning can
help or hurt generalization \citep{amari2020preconditioning}. Our starting
point is full-rank linear regression, where convergence to the unregularized
optimum gives ordinary least squares. The question concerns finite-time
learning and stochastic error; distinct interpolating solutions are a
possible later extension.

Rare predictive features already motivate AdaGrad
\citep{duchi2011adaptive}. More recently, \citet{kunstner2024imbalance}
connect heavy-tailed output-class imbalance to optimizer performance, while
\citet{sreckovic2025batch} show the importance of batch size and careful
tuning in language-model comparisons. Correlated input-feature rarity must
therefore contribute more than another example of slow learning under
imbalance. The candidate contribution is a joint account of dependence,
feature scaling, batch size, and prediction for rare subgroups. The exact
relationship to existing results remains to be established.

The calculations below extend the elementary two-feature flow and scalar
stability results to a fixed common block, finite-update error moments, and
batch-rarity stability under explicit population assumptions. Their relation
to the closest prior theorems remains a research objective. The explicit-penalty
analysis in the appendices supplies complementary baselines, extending the
normalization questions of \citet{larsson2025normalization}. Finite-sample
weighted-lasso distributions and polyhedral selection geometry already have
general treatments \citep{ewald2020distribution,lee2016exact}.

\section{Motivation and setup}

Feature scaling changes the relative penalties in regularized regression.
\citet{larsson2025normalization} study this effect for binary features and
obtain explicit normalization results under orthogonality. The present
investigation asks which conclusions extend to dependent features and which
conditions can replace orthogonality. Our initial focus is on the distinction
between Pearson correlation and covariance relative to a feature's variance.
That distinction becomes consequential when one feature is rare.

Finite-sample distributions for the lasso with correlated designs are already
available \citep{ewald2020distribution}. Polyhedral selection events are also
central to selective inference \citep{lee2016exact}. Accordingly, the candidate
contribution is a characterization of normalization for binary features,
rather than a general distribution theory for correlated lasso estimators.
The references here are a starting point for a targeted novelty review.

Profile out an unpenalized intercept and let $X_c$ and $y_c$ denote the centered
design and response. Write
\[
 C=X_c^\top X_c/n,\qquad z=X_c^\top y_c/n,\qquad
 y_c=X_c\beta^*+\varepsilon_c.
\]
Unless stated otherwise, $C$ is positive definite and inference is conditional
on a fixed design. Independent Gaussian errors with variance
$\sigma_\varepsilon^2>0$ give
\begin{equation}\label{eq:score-law}
 z\sim N(C\beta^*,\tau^2C),\qquad \tau^2=\sigma_\varepsilon^2/n.
\end{equation}
For fixed positive vectors $u,v$, consider the original-scale estimator
\begin{equation}\label{eq:objective}
 \hat\beta=\arg\min_{\beta\in\R^p}
 \left\{\frac12\beta^\top C\beta-z^\top\beta
 +t\sum_j u_j|\beta_j|+\frac r2\sum_j v_j\beta_j^2\right\},
 \qquad t,r\ge0.
\end{equation}
Here $t=\lambda_1/n$ and $r=\lambda_2/n$ correspond to penalties with an
unaveraged squared-error loss. Dividing centered feature $j$ by $s_j$ and
transforming the fitted coefficient back gives $u_j=s_j$ and $v_j=s_j^2$.
This equivalence requires no orthogonality. Allowing separate weights also
covers penalty schemes such as $u_j=v_j=C_{jj}$.

Weights and penalties are fixed with respect to $y$ throughout the distribution
calculations. Cross-validation and response-dependent penalty paths require
separate analysis. Write $\soft_k(x)=\sign(x)(|x|-k)_+$ for
soft-thresholding.

\begin{proposition}[Ridge in arbitrary dimension]\label{prop:ridge}
Let $V=\diag(v)$ and $t=0$. Then
\begin{align*}
 \E\hat\beta&=(C+rV)^{-1}C\beta^*,\\
 \E\hat\beta-\beta^*&=-r(C+rV)^{-1}V\beta^*,\\
 \Cov(\hat\beta)&=\tau^2(C+rV)^{-1}C(C+rV)^{-1}.
\end{align*}
The mean and covariance formulas require only zero-mean errors with covariance
$\sigma_\varepsilon^2 I$, rather than Gaussian errors.
\end{proposition}
\begin{proof}
Differentiate the quadratic objective to obtain $\hat\beta=(C+rV)^{-1}z$,
then take its mean and covariance.
\end{proof}
For $r>0$ and $0<k<1$, the identity $(C+rV)^{-1}C=kI$ implies
$V=(1-k)C/(kr)$. Thus a diagonal ridge penalty can give the same shrinkage
factor for every coefficient vector only when $C$ is diagonal.

\section{Dependence during early stopping}\label{sec:flow}

Gradient flow in the scaled coordinates starts from $\theta(0)=0$ and
satisfies $\dot\theta=S^{-1}(z-CS^{-1}\theta)$. Transforming back gives
\begin{equation}\label{eq:flow}
 \dot\beta=P(z-C\beta),\qquad
 \beta(T)=(I-e^{-TPC})C^{-1}z.
\end{equation}
For the population model, $z=C\beta^*$, so
$\beta(T)=(I-e^{-TPC})\beta^*$.
Although $PC$ need not be symmetric, it is similar to the symmetric positive
definite matrix $S^{-1}CS^{-1}$. Thus the empirical trajectory tends to
$C^{-1}z$ as $T\to\infty$. In this full-rank model scaling changes the path,
while the converged estimator is ordinary least squares.

\begin{proposition}[A transient rare-feature contribution]\label{prop:flow-limit}
Consider population gradient flow from zero for
\[
 C_b=\begin{pmatrix}a&c_b\\c_b&b\end{pmatrix},\qquad
 a>0,\quad b\downarrow0,\quad c_b/b\to\kappa,
\]
with $C_b$ positive definite and fixed $\beta^*\in\R^2$.
Use standard-deviation scaling, so $P_b=\diag(1/a,1/b)$.
For each fixed $T\ge0$,
\begin{align*}
 \beta_{b,1}(T)&\longrightarrow(1-e^{-T})\beta_1^*,\\
 \beta_{b,2}(T)&\longrightarrow(1-e^{-T})\beta_2^*
                         +\kappa T e^{-T}\beta_1^*.
\end{align*}
The convergence is uniform on bounded intervals of $T$.
\end{proposition}
\begin{proof}
The generator satisfies
\[
 P_bC_b\longrightarrow
 L=\begin{pmatrix}1&0\\\kappa&1\end{pmatrix}=I+N,
 \qquad N^2=0.
\]
Continuity of the matrix exponential, uniformly on bounded time intervals,
gives $e^{-TP_bC_b}\to e^{-TL}=e^{-T}(I-TN)$.
Substitute this expression into \eqref{eq:flow} with $z=C_b\beta^*$.
\end{proof}

For binary features with fixed common prevalence $q_1\in(0,1)$ and rare
prevalence $q_2\downarrow0$, write
\[
 a=q_1(1-q_1),\quad b=q_2(1-q_2),\quad c=p_{11}-q_1q_2,
 \quad \max(0,q_1+q_2-1)\le p_{11}\le\min(q_1,q_2).
\]
These bounds imply $c=O(b)$ and hence Pearson correlation tends to zero.
In the feasible nested design $p_{11}=q_2$, however,
$c/b\to1-q_1$. If $\beta_2^*=0$, a nonzero rare coefficient can consequently
appear during training, with limiting value
$(1-q_1)T e^{-T}\beta_1^*$. This is a transient coefficient contribution,
not lasso selection or an asymptotic inconsistency claim. It vanishes as
$T\to\infty$ in the limiting trajectory.

The proposition concerns population dynamics at fixed time. It does not
provide a joint theorem for increasing sample size, changing stopping times,
or stochastic updates. In particular, faster population dynamics do not
establish better performance at a fixed number of processed observations.

\section{Exact moments at a finite update count}\label{sec:discrete}

This section uses the population response model
$y=x^\top\beta^*+\varepsilon$ in a fixed dimension, with known population
centering $\E x=0$ and fixed positive diagonal scales $S$. There is no
intercept update. Assume
\[
 C=\E xx^\top\succ0,\qquad \E\|x\|^4<\infty,\qquad
 \E(\varepsilon\mid x)=0,\qquad
 \E\varepsilon^2<\infty,\qquad
 \E[\varepsilon^2\|x\|^2]<\infty.
\]
The last condition controls gradient noise and does not follow from finite
noise variance alone. It permits heteroscedastic noise and does not require
Gaussian errors or independence of $x$ and $\varepsilon$.
Write $P=S^{-2}$, and use the scaled random feature $z=S^{-1}x$ and
scaled coefficient error
\[
 H=\E zz^\top=S^{-1}CS^{-1},\qquad
 u_k=S(\beta_k-\beta^*),\qquad R=I-\eta H.
\]
Here $z$ denotes a random feature vector, rather than the fixed empirical
score in the setup. Let the initialization have finite second moment and be
independent of all training observations. Deterministic initialization,
including zero, is a special case.

\subsection{Population gradient descent and the online mean}

Population gradient descent on the averaged squared loss obeys
\[
 \beta_{k+1}=\beta_k-\eta PC(\beta_k-\beta^*),\qquad u_{k+1}=Ru_k.
\]
Consequently, for deterministic $\beta_0$ and every integer $k\ge0$,
\begin{equation}\label{eq:discrete-gd}
 \beta_k=\beta^*+(I-\eta PC)^k(\beta_0-\beta^*).
\end{equation}
For random initialization with scaled mean error $m_0$, covariance
$\Sigma_0$, and second moment $U_0$, its mean error is $R^km_0$, its covariance
is $R^k\Sigma_0R^k$, and its second moment is $R^kU_0R^k$.
Thus both matrix moments obey $V_{k+1}=RV_kR$ in population GD.
There is no observation-noise contribution in this
population update. These identities hold for every $\eta>0$ at finite $k$.
Convergence for every initial value holds precisely when
$0<\eta<2/\lambda_{\max}(H)$, since $H$ is symmetric positive definite.

For fresh-batch SGD, let $(x_{ki},\varepsilon_{ki})$, $1\le i\le B$, be
independent copies of $(x,\varepsilon)$, independently across updates and
independently of the initialization. Define
\[
 \widehat H_k=\frac1B\sum_{i=1}^B z_{ki}z_{ki}^\top,\qquad
 \xi_k=\frac1B\sum_{i=1}^B z_{ki}\varepsilon_{ki}.
\]
The averaged stochastic gradient gives
\begin{equation}\label{eq:discrete-sgd-update}
 u_{k+1}=(I-\eta\widehat H_k)u_k+\eta\xi_k,
 \qquad m_k:=\E u_k=R^km_0.
\end{equation}
Thus SGD has exactly the population GD mean at every finite update count.
This identity uses fresh sampling and fixed centering and scaling; it does
not describe repeated passes through a fixed noisy sample. Comparing batch
sizes at a fixed observation budget $N$ requires $k=\lfloor N/B\rfloor$
in these formulas, not a common update count.

\subsection{Second moments and centered covariance}

For symmetric matrices $U$, define
\begin{align}
 \mathcal M(U)&=\E[zz^\top Uzz^\top],\qquad
 \mathcal V(U)=\mathcal M(U)-HUH,
 \qquad Q=\E[\varepsilon^2zz^\top],\label{eq:discrete-operators}\\
 \mathcal T_B(U)&=RUR+\frac{\eta^2}{B}\mathcal V(U)
 =\E[(I-\eta\widehat H_k)U(I-\eta\widehat H_k)].\label{eq:discrete-T}
\end{align}
The identity
$\mathcal V(U)=\E[(zz^\top-H)U(zz^\top-H)]$ shows that $\mathcal V$
maps positive semidefinite matrices to positive semidefinite matrices.
The same is true of $\mathcal T_B$. The fourth moments enter through
$\mathcal M$, including mixed moments between common and rare coordinates.

\begin{proposition}[Finite-update error moments]\label{prop:discrete-moments}
Under the assumptions above, put
$U_k=\E[u_ku_k^\top]$ and
$\Sigma_k=\Cov(u_k)=U_k-m_km_k^\top$.
For every integer $k\ge0$,
\begin{align}
 U_{k+1}&=\mathcal T_B(U_k)+\frac{\eta^2}{B}Q,
 &U_k&=\mathcal T_B^k(U_0)
       +\frac{\eta^2}{B}\sum_{j=0}^{k-1}\mathcal T_B^j(Q),
       \label{eq:discrete-second}\\
 \Sigma_{k+1}&=\mathcal T_B(\Sigma_k)
       +\frac{\eta^2}{B}\{\mathcal V(m_km_k^\top)+Q\}.
       \label{eq:discrete-covariance}
\end{align}
In particular, the centered covariance has the finite-sum decomposition
\begin{equation}\label{eq:discrete-covariance-sum}
 \Sigma_k=\mathcal T_B^k(\Sigma_0)
 +\frac{\eta^2}{B}\sum_{j=0}^{k-1}
    \mathcal T_B^{k-1-j}\{\mathcal V(R^jm_0m_0^\top R^j)\}
 +\frac{\eta^2}{B}\sum_{j=0}^{k-1}\mathcal T_B^j(Q).
\end{equation}
Empty sums are zero, and the zeroth operator power is the identity.
Original-scale quantities are
\begin{equation}\label{eq:discrete-original-moments}
 \E\beta_k-\beta^*=S^{-1}m_k,\qquad
 \Cov(\beta_k)=S^{-1}\Sigma_kS^{-1},\qquad
 \E[(\beta_k-\beta^*)(\beta_k-\beta^*)^\top]=S^{-1}U_kS^{-1}.
\end{equation}
\end{proposition}
\begin{proof}
Conditional on all features in a batch, $\E(\xi_k\mid z_{k1},\ldots,z_{kB})=0$.
Independence of the sampled pairs and the conditional noise-mean assumption
therefore eliminate both cross terms in the second moment of
\eqref{eq:discrete-sgd-update}, even though $\xi_k$ and $\widehat H_k$
need not be independent. They also give
$\E\xi_k\xi_k^\top=Q/B$.
Independence within a batch yields
\[
 \E[\widehat H_kU\widehat H_k]
   =\left(1-\frac1B\right)HUH+\frac1B\mathcal M(U).
\]
Expanding the two update factors proves \eqref{eq:discrete-T}.
Independence from the past then proves the second-moment recurrence.
Subtract $m_{k+1}m_{k+1}^\top=Rm_km_k^\top R$ to obtain
\eqref{eq:discrete-covariance}. Iteration gives both finite sums. The stated
moment assumptions and independence ensure finite moments at every fixed
$k$ by induction. Conjugation by $S^{-1}$ proves the original-scale formulas.
\end{proof}

Equation~\eqref{eq:discrete-covariance-sum} separates propagated initial
uncertainty, multiplicative update variance caused by mean error, and
observation noise. The operator $\mathcal T_B$ also propagates existing
variance through random updates, so the final noise contribution includes
that amplification. With deterministic initialization, $\Sigma_0=0$;
nevertheless $U_0=m_0m_0^\top$ usually is nonzero. In the noiseless case
$Q=0$, SGD may still have positive covariance, whereas initialization at
$\beta^*$ gives no subsequent error. With homoscedastic conditional noise
variance $\sigma_\varepsilon^2$, $Q=\sigma_\varepsilon^2H$.
For an independent test feature from the training population, expected
excess squared loss is
\[
 \frac12\E[(\beta_k-\beta^*)^\top C(\beta_k-\beta^*)]
 =\frac12\operatorname{tr}(HU_k)
 =\frac12m_k^\top Hm_k+\frac12\operatorname{tr}(H\Sigma_k).
\]
These formulas require no stability assumption or inversion of
$I-\mathcal T_B$. For example, column vectorization gives the finite matrix
representation
\[
 \operatorname{vec}\{\mathcal T_B(U)\}=K_B\operatorname{vec}(U),\qquad
 K_B=R\otimes R+\frac{\eta^2}{B}
       \{\E[(zz^\top)\otimes(zz^\top)]-H\otimes H\},
\]
so \eqref{eq:discrete-second} is also an ordinary matrix-power and finite-sum
formula. Mean convergence alone does not ensure bounded second moments.

\subsection{A fixed-update rare-feature mean limit}

The common-block population setup also gives a finite-$k$ limit for both
GD and the mean of fresh-batch SGD. Index populations by $n$, not by the
number of observed training samples, and suppose
\[
 C_n=\begin{pmatrix}A_n&c_n\\c_n^\top&b_n\end{pmatrix}\succ0,\quad
 A_n\to A\succ0,\quad b_n=q_n(1-q_n)\to0,\quad q_n\to0,
 \quad c_n/b_n\to\kappa.
\]
The rare feature is a centered Bernoulli variable. Let
$S_n=\diag(S_{C,n},s_{R,n})$, where $S_{C,n}\to D$ is positive diagonal,
and suppose $b_n/s_{R,n}^2\to\rho\in[0,\infty)$.
Hold $\eta>0$ and $\beta^*$ fixed and take deterministic
$\beta_{n,0}\to\beta^0$ on the original scale. Set
\[
 G=D^{-2}A,\qquad F=I-\eta G,\qquad r_\eta=1-\eta\rho,
 \qquad d=\beta^*-\beta^0.
\]
\begin{proposition}[Rare-feature mean at fixed $k$]\label{prop:discrete-rare}
For each fixed integer $k\ge0$, the original-scale mean coefficients obey
\begin{align}
 \E\beta_{n,C,k}&\longrightarrow\beta_C^*-F^kd_C,\label{eq:discrete-rare-common}\\
 \E\beta_{n,R,k}&\longrightarrow\beta_R^*-r_\eta^kd_R
  +\eta\rho\sum_{j=0}^{k-1}r_\eta^{k-1-j}\kappa^\top F^jd_C.
  \label{eq:discrete-rare-coefficient}
\end{align}
For GD the expectation can be omitted. For SGD the assumptions of
Proposition~\ref{prop:discrete-moments} hold separately in each population;
any finite positive integer batch size $B_n$ gives the same mean.
\end{proposition}
\begin{proof}
The original-scale generator converges to a block triangular matrix:
\[
 P_nC_n\longrightarrow
 L=\begin{pmatrix}G&0\\\rho\kappa^\top&\rho\end{pmatrix}.
\]
For fixed $k$, matrix powering is continuous, so the mean error converges
to $(I-\eta L)^k(\beta^0-\beta^*)$.
The common block follows $F^k$, while solving the rare-coordinate scalar
recurrence driven by the common block gives
\eqref{eq:discrete-rare-coefficient}.
\end{proof}

The term $F^jd_C$ is the remaining common-coefficient error with its sign
reversed. Its convolution with $r_\eta^{k-1-j}$ transfers dependence through
$\kappa$ into the rare mean. For one common feature and standard-deviation
scaling, $G=1$ and $\rho=1$. From zero initialization, for $k\ge1$,
\[
 \E\beta_{n,R,k}\longrightarrow
   \{1-(1-\eta)^k\}\beta_R^*
   +\eta\kappa k(1-\eta)^{k-1}\beta_C^*.
\]
At $\eta=1,k=1$, the zeroth power is one; at $k=0$, the convolution is
zero. For fixed $k$, the general formula also covers $\rho=0$ (including
no rare-feature scaling), when the limiting rare coefficient stays at its
initial value. Variance scaling of the rare coordinate has
$b_n/s_{R,n}^2=1/b_n\to\infty$ and is outside this limit.

Binary feasibility is a distributional requirement, beyond specifying a
positive definite covariance. For common binary vector $W$ and rare
indicator $J$, one can choose distributions of $W\mid J=0$ and $W\mid J=1$
on $\{0,1\}^{p-1}$ and mix them with probabilities $1-q_n$ and $q_n$.
Then $c_n/b_n=\E(W\mid J=1)-\E(W\mid J=0)$ exactly, and nonnegative
conditional cell probabilities ensure feasibility. The numerical checks use
such supports, enumerate independent batches and short update trajectories,
and compare original- and scaled-coordinate moments with the formulas above.

The limit is only for each fixed update count (and hence any fixed finite
set of update counts). It is not uniform over unbounded $k$, does not pass
to an increasing stopping time or a fixed-data sampling theorem, and does
not assert convergence of second moments as $q_n\to0$. In particular,
fourth update moments and noise moments may grow with rarity despite a
finite limiting mean. Taking $\eta\to0$ with $k\eta\to T$ in the limiting
matrix power recovers the block population flow, a separate time-discretization
limit that supplies no joint guarantee for changing $n$, $k$, and $B_n$.

\section{Rare observations and stochastic stability}\label{sec:sgd}

Unit marginal variance does not control the size of individual stochastic
updates. The following scalar calculation isolates this distinction using
known population centering and scaling. It is an elementary moment calculation
that motivates the multivariate problem.

\begin{proposition}[Scalar batch stability]\label{prop:sgd-stability}
Let $X\sim\operatorname{Bernoulli}(q)$ with $0<q<1$, and set
\[
 Z=\frac{X-q}{\sqrt{q(1-q)}},\qquad
 \mu_4=\E Z^4=\frac{(1-q)^3+q^3}{q(1-q)}.
\]
Fit the noiseless model $Y=Z\theta^*$ by SGD on squared loss
$(Z\theta-Y)^2/2$, with constant step size $\eta$ and batches of size
$B\ge1$. Each batch consists of fresh independent observations and is
independent of the previous iterates. For a deterministic initial error
$e_0=\theta_0-\theta^*$,
\[
 \E e_k^2=
 \left[1-2\eta+\eta^2\left\{1+\frac{\mu_4-1}{B}\right\}\right]^k e_0^2.
\]
For $e_0\ne0$, the expected squared error tends to zero precisely when
\begin{equation}\label{eq:sgd-stability}
 0<\eta<\frac{2}{1+(\mu_4-1)/B}.
\end{equation}
\end{proposition}
\begin{proof}
Write $H_k=B^{-1}\sum_{i=1}^B Z_{ki}^2$. The update gives
$e_{k+1}=(1-\eta H_k)e_k$.
Since $\E Z^2=1$, independence within each batch gives
$\E H_k=1$ and $\E H_k^2=1+(\mu_4-1)/B$.
Independence from the past yields the displayed recurrence for the second
moment. Its nonnegative multiplicative factor is less than one exactly
under \eqref{eq:sgd-stability}.
\end{proof}

The population Hessian equals one for every $q$, while $\mu_4\sim q^{-1}$
as $q\downarrow0$. If $B=B_q$ satisfies $B_q q\to0$, the stability bound is
asymptotic to $2B_q q$. If $B_q q\to\infty$, it tends to two, the stability
bound for the population gradient update. At $B_q q\to\gamma\in(0,\infty)$,
it tends to $2\gamma/(1+\gamma)$. These are mean-square stability statements
for the specified online noiseless model, not guarantees for Adam or for
minibatches drawn without replacement from a fixed sample.

This calculation also distinguishes mean convergence from mean-square
convergence: $\E e_k=(1-\eta)^k e_0$, which can converge while the second
moment grows. A covariance-only description of the population loss therefore
does not determine stochastic stability.

In one dimension, rescaling a feature can be absorbed into the learning rate.
The substantive normalization question requires several feature frequencies
and relative coordinate updates. With common and rare predictors together,
a shared learning rate must accommodate their different update moments.
Deriving that interaction with the transient dependence effect in
Section~\ref{sec:flow}, under label noise and a fixed observation budget, is
the next research objective.

\section{A joint observation-budget limit}
\label{sec:budget}

The preceding fixed-update identities do not permit substitution of an
increasing stopping count without further bounds. Here we give such bounds for a
specified last-iterate schedule. The argument uses the classical decreasing-step
framework of \citet{moulines2011nonasymptotic}, with constants controlled as the
fourth update moment grows. It is an existence result under known preprocessing,
with no optimality or novelty claim.

For each sufficiently large integer $n$, let $(W,J) \in \{0,1\}^2$ have genuine
nonnegative joint-cell probabilities, and write
\[
 \begin{aligned}
  p_n & = \E W \to p \in (0,1), & q_n     & = \E J = n^{-1/2}, \\
  a_n & = p_n(1 - p_n),         & b_n     & = q_n(1 - q_n),    \\
  c_n & = \Cov(W,J),            & c_n/b_n & \to \kappa.
 \end{aligned}
\]
Choose $a_0 > 0$ and $n_0$ such that $a_n \ge a_0$ and $q_n \le a_0/64$ for
$n \ge n_0$. Set $x = (W - p_n,J - q_n)^\top$ and
$S_n = \diag(\sqrt{a_n},\sqrt{b_n})$. The response is
$Y = \alpha + x^\top\beta^* + \varepsilon$, with fixed signal $\beta^*$,
$\E(\varepsilon \mid W,J) = 0$, and
$\E(\varepsilon^2 \mid W,J) = \sigma^2 \in (0,\infty)$. The centering, scales,
and intercept $\alpha$ are known and held fixed. Use exactly $n$ iid fresh
observations, zero initialization, and
\begin{equation}
 \label{eq:budget-schedule}
 B_n = 1,\qquad k_n
     = n,\qquad \tau_n
     = \lceil16/q_n\rceil,\qquad \eta_{n,j}
     = \frac{2}{\tau_n + j},\quad 0
     \le j
     < n.
\end{equation}
Thus $n$ is an actual processed-observation budget. All steps and the stopping
count are deterministic, and no iterate averaging is used.

\begin{proposition}[Rare-coefficient error at an observation budget]
 \label{prop:budget}
 Let $u_j = S_n(\beta_{n,j} - \beta^*)$, $m_j = \E u_j$, and decompose
 $\Cov(u_j) = D_j + N_j$, where $D_j$ is the covariance of the noiseless
 trajectory and $N_j$ is the covariance of propagated observation noise. Put
 $C_* = \|S_n\beta^*\|^2$. For every $n \ge n_0$ and $0 \le k \le n$,
 \begin{align}
  |\E\beta_{n,R,k} - \beta_R^*|^2 & \le \frac{C_*}{b_n} \left( \frac{\tau_n}{\tau_n + k} \right)^{7/2},\label{eq:budget-bias} \\
  0 \le (D_k)_{RR}/b_n            & \le \frac{C_*}{b_n} \left( \frac{\tau_n}{\tau_n + k} \right)^3,\label{eq:budget-design}   \\
  0 \le (N_k)_{RR}/b_n            & \le \frac{4\sigma^2}{b_n(\tau_n + k)}.\label{eq:budget-noise}
 \end{align}
 Moreover, $nN_n \to (4\sigma^2/3)I_2$ in operator norm. Consequently,
 \begin{equation}
  \label{eq:budget-mse}
  nq_n\Var(\beta_{n,R,n}) \longrightarrow \frac{4\sigma^2}{3}, \qquad nq_n\E(\beta_{n,R,n} - \beta_R^*)^2
                          \longrightarrow \frac{4\sigma^2}{3}.
 \end{equation}
 The squared bias is $O(n^{-5/4})$, and the original-scale design variance is
 $O(n^{-1})$. These constants are uniform over bounded signals and the stated
 population class with fixed $a_0$ and $\sigma^2$.
\end{proposition}

\begin{proof}
 Suppress $n$ and put $z = S_n^{-1}x$, $H = \E zz^\top$, and
 $K = \E[\|z\|^2zz^\top]$. Binary feasibility gives $|c_n/b_n| \le 1$, so the
 off-diagonal entry of $H$ has magnitude at most $\sqrt{q_n/a_0} \le 1/8$.
 Pointwise, $z_C^2 \le 1/a_0$ and $z_R^2 \le 1/q_n$. Hence
 \[
  \tfrac78 I \preceq H
             \preceq \tfrac98 I,\qquad K
             \preceq \frac{65}{64q_n}H
             \preceq \frac2{q_n}I.
 \]
 For $s_j = \tau_n + j$, let $R_j = I - 2H/s_j$ and
 \[
  \mathcal V(U) = \E[(zz^\top-H)U(zz^\top-H)],\qquad \mathcal T_j(U)
                = R_jUR_j + \frac4{s_j^2}\mathcal V(U).
 \]
 Independence and conditional mean-zero noise give the exact recurrences
 \[
  \begin{split}
   m_{j+1} & = R_jm_j,\quad m_0 = -S_n\beta^*,                                          \\
   D_{j+1} & = \mathcal T_j(D_j) + \frac4{s_j^2}\mathcal V(m_jm_j^\top), \quad D_0 = 0, \\
   N_{j+1} & = \mathcal T_j(N_j) + \frac{4\sigma^2}{s_j^2}H,\quad N_0 = 0.
  \end{split}
 \]
 The two covariance contributions are positive semidefinite. In particular, the
 covariance $D_j$ is distinct from the noiseless second moment
 $U_j^{(0)} = m_jm_j^\top + D_j$, which propagates by $\mathcal T_j$.

 For $U \succeq 0$, the uniform bounds and $q_ns_j \ge 16$ imply
 \[
  \operatorname{tr}\mathcal T_j(U) \le \left( 1 - \frac7{2s_j} + \frac8{q_ns_j^2} \right)\operatorname{tr}U
                                   \le (1 - 3/s_j)\operatorname{tr}U.
 \]
 The product of these factors from $j = l$ to $k - 1$ is at most
 $\{(\tau_n + l)/(\tau_n + k)\}^3$ by an integral bound on the harmonic sum.
 This bounds $\operatorname{tr}U_k^{(0)}$ and therefore $D_k$. Separately,
 $\|R_j\| \le 1 - 7/(4s_j)$ gives the squared-mean exponent $7/2$. Finally,
 $\operatorname{tr}H = 2$ gives
 \[
  \operatorname{tr}N_{j+1} \le (1 - 3/s_j)\operatorname{tr}N_j + 8\sigma^2/s_j^2
                           \le 4\sigma^2/(s_j + 1)
 \]
 by induction from $N_0 = 0$. Dividing the rare-coordinate quantities by $b_n$
 proves the finite bounds.

 To evaluate the noise constant, let $\bar N_0 = 0$ and
 $\bar N_{j+1} = R_j\bar N_jR_j + 4\sigma^2H/s_j^2$. Then
 $F_j = N_j - \bar N_j \succeq 0$ and
 \[
  F_{j+1} = R_jF_jR_j + \frac4{s_j^2}\mathcal V(N_j).
 \]
 Since $\operatorname{tr}\mathcal V(U) \le (2/q_n)\operatorname{tr}U$ for
 $U \succeq 0$, each source has trace at most $32\sigma^2/(q_ns_j^3)$. The
 subsequent contraction product is exactly
 \[
  \prod_{l=j+1}^{k-1}(1 - 3/s_l) = \frac{(s_j - 2)(s_j - 1)s_j}
 {(\tau_n + k - 3)(\tau_n + k - 2)(\tau_n + k - 1)}.
 \]
 Thus
 \begin{equation}
  \label{eq:budget-amplification}
  0 \le \operatorname{tr}F_k \le \frac{32\sigma^2 k}{q_n(\tau_n + k - 3)^3}.
 \end{equation}
 In particular, $n\|F_n\| \to 0$ since $nq_n \to \infty$ and $\tau_n/n \to 0$.

 All $R_j$ commute with $H$, whose two eigenvalues $h$ tend to one. The
 corresponding eigenvalues of $\bar N_n$ are
 \[
  4\sigma^2h\sum_{j=0}^{n-1}s_j^{-2} \prod_{l=j+1}^{n-1}(1 - 2h/s_l)^2.
 \]
 Fix $\delta \in (0,1)$ and $m = \lceil\delta n\rceil$. The terms $j < m$ have
 total trace at most $4\sigma^2(\tau_n + m)^2/(\tau_n + n)^3$, by propagating
 $\bar N_m$ and using the preceding trace bound. After multiplication by $n$,
 their limsup is at most $4\sigma^2\delta^2$. For $m \le j < n$, the uniform
 logarithmic expansion is
 \[
  \log\prod_{l=j+1}^{n-1}(1 - 2h/s_l)^2 = 4h\log\frac{\tau_n + j + 1}{\tau_n + n}
                                          + O((\tau_n + j + 1)^{-1}).
 \]
 The product therefore converges uniformly to $(j/n)^4$ on this late interval.
 Multiplication by $n$ turns its sum into a Riemann sum with limit
 $4\sigma^2\int_\delta^1 x^2\,dx$. First taking $n \to \infty$, then
 $\delta \downarrow 0$, proves $n\bar N_n \to (4\sigma^2/3)I_2$.
 Equation~\eqref{eq:budget-amplification} gives the same limit for $N_n$.
 Finally, $b_n \sim n^{-1/2}$ and $\tau_n \sim 16n^{1/2}$ give the stated bias
 and design rates and \eqref{eq:budget-mse}.
\end{proof}

The result establishes the $(nq_n)^{-1/2}$ second-moment uncertainty scale, with
a schedule-specific constant. It gives neither a central limit theorem nor an
efficiency claim. The same bounds suffice when $q_n \to 0$, $nq_n \to \infty$,
and additionally $n^2q_n^3 \to \infty$; the last condition is sufficient for the
conservative design bound, not claimed necessary. For an auxiliary
iterated-limit comparison, extend the schedule to an infinite fresh-observation
stream. Taking rarity to zero at a fixed update count leaves the original-scale
mean at its zero initialization, while taking the update count to infinity at
fixed full-rank prevalence gives $\beta^*$. The joint path in the proposition
specifies which experiment is analyzed. It does not rank normalization choices
or imply a fixed-data sampling result.

\section{Covariance controls and sampling schemes}
\label{sec:sampling}

Population gradient flow depends on the covariance and fixed scales. SGD also
depends on higher update moments and on how training observations are reused.
This section keeps population centering and scales known and fixes the intercept
at its population value. The fixed-data statements condition on the entire
observed sample, including its realized response noise.

\subsection{Matched covariance, different update moments}

Let $0 < q < 1/2$, $0 \le a \le 1$, and $J \sim \operatorname{Bernoulli}(q)$.
For $U,V \in \{-1,1\}$, specify the conditional joint probabilities by
\[
 \Prob(U = u,V = v \mid J = j) = \frac{1 + a_j uv}{4},\qquad a_1
                               = a,\quad a_0
                               = -\frac{qa}{1 - q}.
\]
All cells are nonnegative. The original binary features are $W_1 = (U + 1)/2$,
$W_2 = (V + 1)/2$, and $J$. Their means and covariance are $(1/2,1/2,q)$ and
$\diag(1/4,1/4,q(1 - q))$ for every $a$. After standard-deviation scaling, put
$z = (U,V,Z)^\top$, $Z = (J - q)/\sqrt{q(1 - q)}$. Then
\[
 H = I_3,\qquad \E[UVZ^2]
   = d_a
   := \frac{a(1 - 2q)}{1 - q},\qquad \E Z^4
   = \mu_4.
\]
Thus even all marginal feature distributions agree, while a mixed fourth moment
changes.

\begin{proposition}[A covariance-matched SGD comparison]
 \label{prop:matched}
 For the distributions above, population GD, population gradient flow, and the
 fresh-SGD mean coincide under identical scales and initialization. In the
 noiseless model, initialize with original-scale coefficient error
 $(0,0,e_0)^\top$. At the same $B$ and $\eta$, after two fresh-batch updates,
 \[
  \E_a[(\beta_{R,2} - \beta_R^*)^2] - \E_0[(\beta_{R,2} - \beta_R^*)^2] = \frac{2\eta^4d_a^2}{B^2}e_0^2.
 \]
 The comparison uses $2B$ processed observations in both populations.
\end{proposition}
\begin{proof}
 Conditional balanced signs give zero means and zero common--rare covariances.
 The choice $(1 - q)a_0 + qa_1 = 0$ gives $\E UV = 0$, while conditioning on $J$
 gives the stated $d_a$. The mean claims follow from $H = I_3$. Write
 $E_{33} = e_3e_3^\top$. Then
 \[
  \mathcal M_a(E_{33}) = \begin{pmatrix}
                          1   & d_a & 0     \\
                          d_a & 1   & 0     \\
                          0   & 0   & \mu_4
                         \end{pmatrix}.
 \]
 Since $R = (1 - \eta)I_3$, the only $a$-dependent entries of
 $\mathcal T_{a,B}(E_{33})$ are its two common off-diagonal entries
 $\eta^2d_a/B$. Applying $\mathcal T_{a,B}$ again contributes $2\eta^4d_a^2/B^2$
 to its rare diagonal entry. Multiplication by the initial scaled squared error
 $q(1 - q)e_0^2$ and division by $q(1 - q)$ prove the result.
\end{proof}

Stability can differ as well. At $q = 1/4$, $\mu_4 = 7/3$ and $d_a = 2a/3$. On
symmetric matrices, the eigenvalues of $\mathcal M_a$ are
\[
 0,\quad 2,\quad 2(1 - d_a),\quad 2(1 + d_a),\quad \frac{2 + \mu_4 \pm \sqrt{(\mu_4 - 2)^2 + 8(1 + d_a^2)}}{2}.
\]
These follow by separating the common diagonal difference, the two common--rare
off-diagonal modes, and the span of the common diagonal sum, the common
off-diagonal, and the rare diagonal. The largest eigenvalue $\lambda_a$ is the
displayed plus root, so
\[
 \eta_{\mathrm{crit},a} = \frac{2}{1 + (\lambda_a - 1)/B}.
\]
For $B = 1$ the thresholds are $12/(13 + \sqrt{73})$ at $a = 0$ and
$12/(13 + \sqrt{105})$ at $a = 1$. The common step $\eta = 8/15$ is strictly
stable for the first distribution and fails all-initializations mean-square
convergence for the second. Both means converge. This comparison isolates
information absent from the population covariance. For fixed $B$ and
$q \downarrow 0$, both thresholds remain asymptotic to $2Bq$, consistent with
the leading rare-limit result in Section~\ref{sec:sgd}.

\subsection{Sampling a fixed dataset with replacement}

For fixed observations $(z_i,y_i)_{i=1}^n$, suppose
$H_D = n^{-1}\sum_i z_i z_i^\top \succ 0$, and let $\widehat\theta_D$ be the
empirical least-squares coefficient in scaled coordinates. Define
\[
 r_i = y_i - z_i^\top\widehat\theta_D,\quad g_i
     = z_i r_i,\quad \frac1n\sum_i g_i
     = 0,\quad v_k
     = \theta_k - \widehat\theta_D,\quad R_D
     = I - \eta H_D.
\]
Each update samples $B$ indices independently and uniformly with replacement,
independently of previous index draws conditional on the dataset. Write
\[
 \mathcal M_D(U) = \frac1n\sum_i z_i z_i^\top U z_i z_i^\top,\quad \mathcal V_D(U)
                 = \mathcal M_D(U) - H_DUH_D,\quad Q_D
                 = \frac1n\sum_i g_i g_i^\top,
\]
and let $\mathcal T_{D,B}(U) = R_DUR_D + \eta^2\mathcal V_D(U)/B$. The
additional linear map
\[
 J_D(m) = \frac1n\sum_i z_i z_i^\top m g_i^\top
\]
records correlation between the sampled curvature and the residual score.

\begin{proposition}[Conditional replacement moments]
 \label{prop:replacement}
 Conditional on the dataset, let $m_k = \E_Dv_k$ and $U_k = \E_D[v_kv_k^\top]$,
 with finite initial second moment and initialization independent of the sampled
 indices. Then
 \begin{align*}
  m_{k+1} & = R_Dm_k,                                                                     \\
  U_{k+1} & = \mathcal T_{D,B}(U_k) + \frac{\eta^2}{B}\{Q_D - J_D(m_k) - J_D(m_k)^\top\}.
 \end{align*}
 For $\Sigma_k = U_k - m_km_k^\top$, the covariance recurrence has the
 additional term $\eta^2\mathcal V_D(m_km_k^\top)/B$ on the right.
\end{proposition}
\begin{proof}
 Put $\widehat H = B^{-1}\sum_b z_{i_b}z_{i_b}^\top$ and
 $\bar g = B^{-1}\sum_b g_{i_b}$. The update is
 $v_{k+1} = (I - \eta\widehat H)v_k + \eta\bar g$. The normal equations give
 $\E_D\bar g = 0$ and $\E_D\bar g\bar g^\top = Q_D/B$. Independence within a
 batch and $n^{-1}\sum_i g_i = 0$ give
 $\E_D[\widehat Hm\bar g^\top] = J_D(m)/B$. Expanding the update outer product
 proves the second-moment formula. The conditional mean and covariance follow
 directly.
\end{proof}

Thus empirical normal equations do not replace the fresh-sample conditional
noise-mean assumption. The $J_D$ terms can be nonzero: scalar observations
$z_1 = 1,z_2 = 2$ with residuals $r_1 = 2,r_2 = -1$ have zero empirical score
but $J_D(m) = -3m$. The empirical mean of these features need not vanish when
centering uses known population quantities. The conditional mean follows
empirical GD toward $\widehat\theta_D$, and the original-scale error relative to
the population signal is
\[
 \beta_k - \beta^* = S^{-1}v_k + (S^{-1}\widehat\theta_D - \beta^*).
\]
The second term persists when the same noisy observations are reused. In the
noiseless full-rank model the residuals vanish, so the homogeneous operator
criterion in Proposition~\ref{prop:matrix-stability} applies with empirical
moments. Population thresholds involving $Bq$ require additional empirical
moment control before being used conditionally on a random dataset.

\subsection{Random reshuffling}

For single-observation updates, augment the error by a constant coordinate:
\[
 w = \begin{pmatrix}
      v \\
      1
     \end{pmatrix},\qquad A_i
   = \begin{pmatrix}
      I - \eta z_i z_i^\top & \eta g_i \\
      0                     & 1
     \end{pmatrix}.
\]
For a uniform random permutation $\pi$ of the $n$ indices, put
$A_\pi = A_{\pi(n)}\cdots A_{\pi(1)}$. Independently reshuffling each epoch
gives exact conditional epoch recurrences
\[
 \E_Dw_{e+1} = \E_\pi[A_\pi]\E_Dw_e,\qquad \E_D[w_{e+1}w_{e+1}^\top]
             = \E_\pi\left[ A_\pi\E_D[w_ew_e^\top]A_\pi^\top \right].
\]
Here $e$ counts epochs, each consisting of $n$ updates and $n$ processed
observations. These identities follow by conditioning on the epoch's starting
iterate and using an independent permutation. They do not factor into $n$
independent single-update operators. Fixed-size disjoint minibatches are handled
by replacing each $A_i$ with its batch-average update, specifying any incomplete
final batch.

If the original dataset consists of independent identically distributed
observations, an independent permutation preserves its product distribution.
Consequently, the first pass, unconditionally over the dataset, has exactly the
fresh-sample law for an initialization independent of the dataset and known
population centering and scales. Conditioning on the dataset or making another
pass changes this conclusion. Empirical centering, estimated scales, and an
initialization fitted to the dataset also invalidate this argument.

For a concrete conditional counterexample, take a population-standardized
Bernoulli feature with prevalence $1/3$ and a noiseless dataset containing one
observation at each of its values $\sqrt2$ and $-1/\sqrt2$. The two curvatures
are $2$ and $1/2$. At a common two-observation budget,
\[
 v_2^{\mathrm{shuffle}} = (1 - 2\eta)(1 - \eta/2)v_0,\qquad \E_Dv_2^{\mathrm{replace}}
                        = (1 - 5\eta/4)^2v_0.
\]
The difference is $-9\eta^2v_0/16$. Replacement contracts in mean square exactly
for $0 < \eta < 20/17$, while reshuffling contracts per epoch for
$0 < \eta < 5/2$. At $\eta = 2$, reshuffling reaches zero in one epoch and
replacement has second-moment factor $9/2$ per update. These statements follow
by multiplying the two scalar factors and averaging their squares, respectively.
They concern this specified dataset and do not establish a general advantage for
reshuffling.

Reusing response noise also changes unconditional error. With one observed pair,
an equiprobable sign feature $z$, independent mean-zero noise of variance
$\sigma^2$, and initialization at truth, repeated use gives
\[
 u_k = \{1 - (1 - \eta)^k\}z\varepsilon,\qquad \E u_k^2
     = \sigma^2\{1 - (1 - \eta)^k\}^2.
\]
Fresh observations instead give $\eta^2\sigma^2\sum_{j=0}^{k-1}(1 - \eta)^{2j}$.
For $0 < \eta < 2$, the limits are $\sigma^2$ and $\eta\sigma^2/(2 - \eta)$,
respectively. Conditional finite-data identities also require extra
integrability before averaging over a random dataset: noiseless reuse of a
single feature gives $u_2^2 = (1 - \eta z^2)^4u_0^2$, which can have infinite
expectation when $\E z^4 < \infty$ but $\E z^8 = \infty$.

\section{Controlled rare-subgroup prediction pilot}
\label{sec:pilot}

We compare four fixed scalings on six feasible two-binary-feature designs. Each
design has $n = 400$ observations, common prevalence $p = 1/2$, and rare
prevalence $q \in \{0.05,0.01\}$. The joint law is independent, nested ($J = 1$
implies $W = 1$), or negatively associated at the disjoint boundary
$\Prob(W = 1,J = 1) = 0$. Training cell counts equal their population
expectations, giving 20 or 4 rare observations. The signals are $(1,0)$,
$(0,1)$, $(1,1)$, and $(1,-1)$, with iid Gaussian observation noise of variance
one. Population centering and scales are known, the centered-model intercept is
fixed at zero, and coefficients start at zero. Repeated visits to a training
observation reuse its noise. This experiment concerns fixed-data sampling,
separate from the fresh-sample joint limit in Section~\ref{sec:budget}.

For coordinate variance $v_j$, the feature divisors are $1$, $\sqrt{v_j}$,
$v_j$, and $\sqrt{v_j + 0.01}$. We use full-batch GD and replacement SGD with
$B \in \{1,20\}$. Budgets of $M \in \{400,2000,8000\}$ processed observations
give $M/400$ GD updates and $M/B$ SGD updates. Every design, signal, method,
scaling, and budget receives separate constant-step tuning over the same nine
factors $a \in \{0.001,0.003,0.01,0.03,0.1,0.3,0.7,1.5,2.5\}$, with
$\eta = a/\lambda_{\max}(S^{-1}CS^{-1})$. The objective is mean rare-subgroup
risk across 32 tuning noise realizations. Evaluation uses 256 disjoint noise
realizations, paired across methods and scalings. This is oracle tuning with the
true signal and test law, rather than validation-based tuning. Nominal tuning
effort is $9 \cdot 32 \cdot M$ processed observations per comparison.

For original-scale error $e = \beta - \beta^*$, let $V_D = \E[ee^\top\mid D]$,
where $D$ is the fixed noisy training dataset and the expectation integrates
replacement-index sampling. We report excess squared prediction errors
\[
 R_D = \operatorname{tr}(CV_D),\qquad R_{D,R}
     = \operatorname{tr}(C_RV_D),\qquad C_R
     = \E[xx^\top\mid J=1].
\]
These reported risks omit the factor $1/2$ in the training loss and exclude
irreducible test noise. The conditional second moment $C_R$ includes the nonzero
conditional feature mean. An isolated rare-coefficient error of one therefore
contributes $q(1 - q)$ to aggregate risk and $(1 - q)^2$ to rare risk, a ratio
of 99 at $q = 0.01$.

Conditional means and second moments are computed exactly by powering the affine
replacement operator of Section~\ref{sec:sampling}, augmented by a constant
coordinate. This retains the empirical score and curvature--noise cross terms.
Independent enumeration of ordered batches and two-update paths, direct
prediction-risk calculations, and a separate moment recurrence check agree to
numerical precision. Monte Carlo standard errors (MCSEs) below concern training
noise only, conditional on the fixed designs and selected steps. They exclude
variation in tuning, random designs, and estimated preprocessing.

\begin{table}[t]
 \centering
 \small
 \setlength{\tabcolsep}{4pt}
 \begin{tabular}{lrllll}
  \toprule
  Method        &  $M$ & None          & Std. dev.     & Variance      & Damped        \\
  \midrule
  GD            &  400 & .2954 (.0013) & .1565 (.0026) & .1948 (.0024) & .1704 (.0026) \\
  GD            & 2000 & .2028 (.0023) & .1140 (.0031) & .1415 (.0037) & .1161 (.0035) \\
  GD            & 8000 & .1022 (.0022) & .1152 (.0034) & .1308 (.0046) & .1190 (.0033) \\
  SGD, $B = 1$  &  400 & .2275 (.0024) & .2119 (.0027) & .2554 (.0028) & .2228 (.0018) \\
  SGD, $B = 1$  & 2000 & .1384 (.0035) & .1565 (.0045) & .1914 (.0045) & .1506 (.0033) \\
  SGD, $B = 1$  & 8000 & .1257 (.0036) & .1795 (.0060) & .1995 (.0059) & .1662 (.0058) \\
  SGD, $B = 20$ &  400 & .2145 (.0017) & .2213 (.0023) & .2594 (.0022) & .1916 (.0024) \\
  SGD, $B = 20$ & 2000 & .1265 (.0022) & .1513 (.0035) & .1798 (.0030) & .1282 (.0029) \\
  SGD, $B = 20$ & 8000 & .1058 (.0030) & .1290 (.0029) & .1522 (.0033) & .1180 (.0037) \\
  \bottomrule
 \end{tabular}
 \caption{
  Rare-presence excess squared prediction error, with MCSE in parentheses. Each
  entry averages equally over all six designs and four signals. MCSEs preserve
  the pairing across signals within a design. Equal weighting is a descriptive
  choice for this finite grid.
 }
 \label{tab:pilot}
\end{table}

Table~\ref{tab:pilot} shows substantial differences and changes in their
direction. At $q = 0.01$ with independent features and $M = 400$,
standard-deviation scaling changes rare risk relative to no scaling by
$-0.58017$ (paired MCSE $0.01517$) for the rare-only signal and $+0.14895$
($0.01301$) for the common-only signal. This budget permits only one GD update.
Dependence also reverses equal-sign and opposite-sign comparisons. Across all 24
cases, the average GD advantage of standard-deviation scaling at $M = 400$
reverses by $M = 8000$. The corresponding reversal for $B = 1$ SGD occurs by
$M = 2000$. These are examples of finite-budget preconditioning and
noise-fitting tradeoffs, without a general ordering of the scaling rules.

The finite tuning grid is a material limitation. Its smallest factor is selected
in 195 of 864 evaluated settings and its largest in nine. All selected endpoints
are finite, but nine GD selections are asymptotically unstable and are evaluated
after only one or five updates. Of 7776 candidate endpoints, 2160 have
homogeneous second-moment spectral radius at least one and 1028 have nonfinite
or invalid moments; these categories overlap. All outcomes are retained. The
grid was fixed before evaluation and was not expanded after observing these
boundary selections. The repository's \texttt{notes/rare-subgroup-pilot.md}
records the full protocol, aggregate risks, casewise contrasts, limitations, and
reproduction commands. The generated JSON retains coefficients, covariance,
risks, counts, scales, selected rates, seeds, and software provenance. The
experiment supplies a controlled illustration, without establishing robustness
to a wider tuning search, estimated preprocessing, or real-data applications.

\section{Estimated preprocessing and validation selection}
\label{sec:practical}

We next examine whether the effects in the controlled comparisons remain visible
when we sample training data, estimate preprocessing, and select configurations
by validation. We use a finite grid of candidates, so this experiment does not
extend the population theorems to estimated scales. The full protocol, results,
and reproducible artifacts are recorded in \texttt{notes/practical-selection.md}
and \texttt{notes/practical-selection-results.md}.

For the synthetic experiment, we use eight feasible binary laws with 128
independent replications per law, 256 training observations, and 128 independent
validation observations. We include independent and nested signals at $q = 0.02$
and $q = 0.002$, signals with opposite signs under nested and disjoint
dependence, and rare-only and common-only signals. We compare known
preprocessing with estimated scales, training centering with a fixed or fitted
intercept, and uncentered predictors with a fitted intercept. We estimate
variances with denominator $n$ and give constant training columns divisor one
and coefficient zero. Absent features therefore cannot acquire a fitted effect
in the estimated preprocessing arms. Fitted intercepts start at the training
response mean, while the population-oracle arms start at the true centered
intercept of zero. SGD samples with replacement from a fixed noisy training
set, so revisiting an observation also reuses its noise.

The divisors are $1$, $\sqrt{v_j}$, and $\sqrt{v_j + 0.01}$. For each scaling
and preprocessing arm, we use five rates $\eta = a/\max(L,10^{-12})$,
$a \in \{0.001,0.01,0.05,0.2,0.8\}$, and two batch sizes. Here $L$ is the
largest empirical curvature, including the unscaled intercept coordinate when
fitted. This calibration does not guarantee stochastic stability. We use batches
$B \in \{1,16\}$ and observation caps $M \in \{128,512,2048\}$. We retain
initialization and every checkpoint, and compare fixed $a = 0.05$, $B = 16$ with
validation selection of the rate and batch size, or of the rate, batch size,
and stopping count. We select by either aggregate or rare-presence validation
half squared loss. When the rare validation group is empty, its selection rule
falls back to aggregate loss. We use no test outcomes in selection and do not
refit selected models on validation data. Each combination of scaling and
preprocessing uses ten trajectories, costing $10M$ processed observations at
cap $M$, with checkpoints sharing paths.

Differences in how quickly the coefficients are learned remain visible with
estimated preprocessing. For the independent $(1,1)$ signal at $q = 0.02$ and
noise variance three, the fixed configuration gives an SD-minus-none rare excess
half squared loss of $-0.1653$ (paired MCSE $0.0295$) at $M = 128$ and
$+0.0901$ ($0.0697$) at $M = 2048$ with training centering and a fitted
intercept. The early advantage is clear, while the later disadvantage is within
two MCSE of zero. At the latter budget, the rare-only signal favors SD by
$0.2320$ ($0.0233$), whereas the common-only signal favors no scaling by
$0.1272$ ($0.0159$). These observations concern the declared configurations with
sampled data, rather than the exact calibrations of the population comparisons.

Missing groups create a further problem. At $q = 0.002$, the two anchor laws
have no rare validation observations in 106 and 105 of 128 replications.
Validation can then select a finite but explosive rare coefficient because no
rare cases are available to reveal its error. One SD fit with centered
predictors has a rare coefficient of about $-6.719 \times 10^{16}$ and
validation half squared loss $2.086$, yet its exact rare loss over the binary
cells is about $2.257 \times 10^{33}$. We retain this path in the results. In
a separate investigation, we replay every eligible trajectory for the worst
selected replication in each of 24 combinations of law and scaling, using
160-digit decimal arithmetic. All 24 selections remain unchanged, so the
detected floating-point cancellation does not explain the failure. Means
dominated by these extreme finite outcomes do not precisely estimate population
risk. We also retain all 151 failed synthetic endpoints, flagged in the archive.

\paragraph{A sparse categorical task.}

We fit a squared-loss linear predictor to the recorded poisonous label in the
UCI Mushroom dataset \citep{uciMushroom}, using all 22 categorical predictors.
The declared subgroup consists of mushrooms with green spore-print color, which
accounts for 72 of 8124 records. The records are hypothetical descriptions
derived from a field guide, so evaluation on random partitions does not
establish generalization to new species or field data. There are no duplicate
predictor rows. We construct the vocabulary and estimate centering and scales
from training data alone, treat missingness as a category, and map unknown
categories to zero blocks. We store the one-hot design sparsely and implement
centering algebraically. Because the full one-hot design is rank deficient, its
coefficients depend on the optimizer and cannot be interpreted as uniquely
identified effects.

We use 32 random partitions for each pair of training and validation sizes,
$(256,128)$ and $(2048,1024)$, holding out the remaining records for testing.
All candidates share partitions and replacement index streams. We use batches
$B \in \{16,128\}$ and caps $M \in \{512,2048,8192\}$, and do not clip
predictions. Table~\ref{tab:practical} reports results with centered predictors
and a fitted intercept at the largest cap, jointly selecting the rate, batch
size, and stopping count. We separately retain the other preprocessing arms,
shorter caps, fixed stops, all paths, and the search across scaling rules, which
requires three times as much work.

\begin{table}[t]
  \centering
  \small
  \begin{tabular}{llrr}
    \toprule
    Validation objective & Scaling &    Aggregate loss &         Rare loss \\
    \midrule
    Aggregate            & None    & .001754 (.000060) & .012482 (.001386) \\
    Aggregate            & SD      & .004617 (.002125) & .005912 (.001054) \\
    Aggregate            & Damped  & .001175 (.000034) & .004151 (.000465) \\
    Rare                 & None    & .001878 (.000130) & .012465 (.001382) \\
    Rare                 & SD      & .123905 (.075527) & .003158 (.000451) \\
    Rare                 & Damped  & .014966 (.009145) & .003984 (.000495) \\
    \bottomrule
  \end{tabular}
  \caption{
    Mushroom held-out half squared loss (MCSE), with $n = 2048$, 1024 validation
    records, and budget cap $M = 8192$. Each scaling receives the same ten
    trajectories. MCSE describes variation across partitions and optimization
    runs conditional on this dataset. It does not give a population confidence
    interval.
  }
  \label{tab:practical}
\end{table}

Under aggregate selection, the SD-minus-none rare-loss difference is $-0.00657$
(paired MCSE $0.00153$), and the aggregate difference is $+0.002863$
($0.00213$). Selecting SD by rare rather than aggregate validation loss reduces
rare loss by $0.002754$ ($0.000967$), while increasing aggregate loss by
$0.11929$ ($0.0756$). The green coefficient on the original scale averages
$0.4822$, $0.7462$, and $0.7678$ for none, SD, and damped scaling under
aggregate selection. Other categorical coefficients partly offset its larger
contribution. These results show that relative scaling affects prediction under
a finite budget, with the ordering depending on the objective and tuning.

At the smaller sample size, five training partitions and thirteen validation
partitions contain no green observations. At the larger size, aggregate
validation selects the largest rate factor for unscaled training in every
partition, leaving the optimum unresolved. These limits in coverage and search
range preclude a general normalization recommendation. Independent checks
reconstruct every archived validation and test loss, selection, and paired
uncertainty estimate, and replay representative complete paths in transformed
coordinates. We preserve the earlier controlled studies and their reported
failures.

\section{Variance scaling across practical preprocessing arms}
\label{sec:variance-scaling-followup}

In a separately declared follow-up, we add variance scaling, $s_j = v_j$
($\delta = 1$), to every preprocessing arm of the sampled binary and Mushroom
studies. We reuse all eight binary laws with 128 replications per law, both
Mushroom sample sizes with 32 partitions per size, and the original training,
validation, and test samples. We use population variances only in the synthetic
arm with known preprocessing. All other arms use training variances with
denominator $n$. We keep the same treatment of constant columns and unknown
categories, as well as the same centering, initialization, and intercept rules.

We give no scaling, SD scaling, damping $.01$, and variance scaling the same
25 rate factors relative to empirical curvature. These comprise the earlier
17-factor Mushroom grid, two smaller factors $10^{-5}$ and $10^{-4}$, and six
successive doublings from $12.8$ through $409.6$. With two batch sizes, each
rule has nominal cost $50M$ processed observations per preprocessing arm and
replication or partition at cap $M$. We jointly select the rate, batch size, and stopping
count by aggregate or rare validation loss. The primary search uses the original
checkpoints, while a secondary search adds intermediate checkpoints. We also
retain the nested five-factor practical and 17-factor Mushroom searches to
compare results at the earlier tuning effort. Test outcomes enter neither
selection nor the choice of grid. We retain initialization, fallbacks for empty
groups, numerical failures, and extreme finite selections.

The motivation for variance scaling under the lasso concerns penalty weights on
the original coefficient scale. Here, we use no explicit penalty. The SGD
multiplier for an active coordinate is $\eta/s_j^2 = \eta/v_j^2$, and
calibrating by empirical curvature also changes $\eta$. We therefore compare
how the rules affect relative learning, noise, and stability under a finite
budget. The motivation from shrinkage does not establish which rule should have
lower prediction risk after tuning.

\paragraph{Mushroom results and search boundaries.}

At $n = 2048$ and $M = 8192$, aggregate selection with the original stops
gives variance scaling aggregate loss $.311988$ (MCSE $.201772$) and rare
loss $.051586$ ($.008319$) when predictors are centered and the intercept is
fitted. Its rare-loss differences relative to none, SD, and damping are
$.045630$ ($.008323$), $.046822$ ($.008189$), and $.050904$ ($.008322$),
respectively. Each exceeds two paired MCSE, while the aggregate differences do
not. For variance scaling, aggregate loss has median $.069802$ and maximum
$6.505207$. Selecting by rare validation raises its aggregate mean to $25.5720$
(MCSE $8.0453$).

For the two primary selections at factor $6.4$, the newly tested neighbor
$12.8$ gives finite but worse validation scores. Across both Mushroom sample
sizes, all preprocessing arms, caps, objectives, and both expanded stopping
searches, no selection with a positive budget reaches either grid boundary.
With the original stops, every selected predictor matches the nested 17-factor
search. Sixty candidate indices change because the deterministic tie rule
assigns the same initialization a lower factor. The refined stopping search
selects some added rates, up to $102.4$, and we report it separately.

\begin{table}[ht]
  \centering
  \begin{tabular}{rlrrrr}
    \toprule
     $n$ & Preprocessing      &    None &      SD &  Damped & Variance \\
    \midrule
     256 & Centered, fixed    & .101062 & .075937 & .077910 &  .079580 \\
     256 & Centered, fitted   & .098799 & .076035 & .078161 &  .079579 \\
     256 & Uncentered, fitted & .136151 & .076391 & .082727 &  .079542 \\
    2048 & Centered, fixed    & .005647 & .004722 & .000652 &  .051601 \\
    2048 & Centered, fitted   & .005955 & .004764 & .000682 &  .051586 \\
    2048 & Uncentered, fitted & .043284 & .005060 & .007037 &  .049465 \\
    \bottomrule
  \end{tabular}
  \caption{
    Mean Mushroom rare-group loss under aggregate validation, cap $8192$, and
    the original stops, using the same 25-factor search for every rule. The
    artifacts report both risks, paired uncertainty, and extreme outcomes for
    every preprocessing arm.
  }
\end{table}

At $n = 256$ with centered predictors and a fitted intercept, variance scaling
increases aggregate loss relative to none by $.007047$ (MCSE $.000555$). The
rare-loss difference of $-.019220$ ($.012351$) is unresolved. Five training
partitions and thirteen validation partitions lack green observations, and we
retain all of them. With uncentered predictors, variance scaling gives aggregate
means $.179701$ (MCSE $.150890$) and $13.787223$ ($12.294085$) at the two
sample sizes. These MCSEs describe variation across partitions of one fixed,
previously studied dataset.

\paragraph{Extreme synthetic outcomes and unresolved cases.}

We evaluate all eight laws and five preprocessing arms at equal effort. For the
rarer C2 law with centered predictors and a fitted intercept, aggregate
selection at cap $2048$ gives variance scaling an exact rare-risk median of
$.6836$, mean $1.727 \times 10^{30}$, MCSE $1.552 \times 10^{30}$, and
maximum $1.974 \times 10^{32}$. Extreme outcomes dominate these means,
including those for SD and damping, so they do not precisely estimate expected
risk. Empty rare validation groups can still conceal severe prediction errors.
For the rare-only signal, variance scaling gives aggregate mean $.01065$
($.00311$) but rare mean $.4122$ ($.1467$), compared with $.2326$ ($.0319$)
for no scaling.

Across the two expanded stopping searches, synthetic selections reach the lower
rate boundary 106 times and the upper boundary 52 times, counting repeated
queries on the same replications. Only variance scaling reaches the upper
boundary, under known preprocessing in the two rarer anchor laws. Resolving the
Mushroom boundaries therefore leaves these synthetic cases open. Across both
stages, the archives retain 6,451,200 endpoints and 387,072 choices, including
1,031,540 failed endpoints on 330,918 failed paths.

\paragraph{Independent audit and numerical limits.}

The audit reconstructs every loss and selection, verifies all 179 preserved
files, and replays 100,271 paths in explicitly transformed coordinates.
Overlapping frozen coefficient paths reproduce exactly. Reconstructed losses
have maximum normwise discrepancy $1.36 \times 10^{-14}$ and leave every
selection unchanged. Mushroom passes the declared trajectory and prediction
tolerances. For synthetic data, however, the audit retains 63 discrepancies in
twelve replications, all of which have empty rare validation groups. These
comprise 60 checks of selected predictions, two first-failure times, and one
trajectory check. Repeated queries can refer to the same selected path. The
maximum normwise coefficient and prediction discrepancies are $.003747$ and
$.004184$, exceeding the declared tolerances of $2 \times 10^{-9}$ and
$10^{-8}$. All stored checkpoint statuses agree.

The transformed replay therefore does not certify these extreme synthetic
paths. We retain the original candidates and selections in the archive. The
result note records the affected configurations and explains how we extended the
checker to retain disagreements in failure times rather than stop at the first
one. See \path{notes/variance-scaling-followup-results.md} for the protocol,
archives, and reproduction commands. All earlier studies retain their frozen
file hashes.

\section{Research program and evaluation}

The first milestone is a finite-time theory for one rare binary feature and a
fixed block of common predictors, supported by a controlled experiment and a
comparison with the closest prior theorems. Write
\[
 C_n = \begin{pmatrix}
        A_n      & c_n \\
        c_n^\top & b_n
       \end{pmatrix},\qquad A_n
     \to A
     \succ 0,\quad b_n
     \downarrow 0,\quad c_n/b_n
     \to \kappa.
\]
Sections~\ref{sec:flow}--\ref{sec:sgd} establish the population block-flow
limit, exact finite-update GD and fresh-SGD error moments, and stability
thresholds for a bounded common block across the three $Bq$ regimes. The mean
trajectory separates incomplete common-feature learning from direct rare-feature
learning, while the covariance formulas separate initial uncertainty,
multiplicative update variance, and observation noise. These results use known
population centering and fixed scales. Section~\ref{sec:sampling} supplies a
matched-covariance binary comparison, conditional fixed-data replacement
identities, and exact reshuffling epoch recurrences. Counterexamples delimit
their relation to fresh sampling. Section~\ref{sec:budget} supplies a joint
observation-budget limit for a specified decreasing-step last iterate with one
common and one rare binary feature. Section~\ref{sec:pilot} supplies a
controlled fixed-data comparison with consequential rare-subgroup risk
differences and reversals. Broader joint sampling and stopping-time limits
remain open.

For empirical samples, specify $n \to \infty$, rare prevalence
$q_n \downarrow 0$, and whether $nq_n \to \infty$ or rare counts remain bounded.
Compare error contributions with statistical uncertainty, typically of order
$(nq_n)^{-1/2}$ under suitable regularity conditions. Analyze the effects of
estimating centering and scales, and distinguish a stopping time fixed in
advance from one chosen by validation. Full-rank convergence to OLS,
early-stopping regularization, and selection among interpolating solutions are
separate questions.

The initial experiment holds binary cell counts and population preprocessing
fixed. Section~\ref{sec:practical} adds sampled designs, estimated
preprocessing, intercept ablations, and validation selection, followed by a
sparse categorical benchmark. It retains effects of relative learning but also
exposes missing-group failures and rate-grid boundaries. Wider rate ranges and
group-coverage-aware selection are hypotheses for a further, separately declared
study. AdaGrad, Adam, or a neural model should be added only to test a specific
mechanism, with momentum, stabilizing constants, clipping, schedules, and weight
decay recorded. A fixed diagonal preconditioner is not a model of their complete
dynamics.

Report performance against processed observations as well as update count;
changing batch size changes the cost of an update. Measure original-scale
coefficient error, aggregate excess prediction risk, and excess risk conditional
on rare-feature presence. An isolated rare-coefficient error of constant size
contributes only $O(q_n)$ to population risk. To assess a consequence beyond
aggregate averages, also consider a specified test prevalence shift that
preserves the conditional response model and uses the training centering and
scales at prediction time. New intercept or centering errors must not be
confused with coefficient effects.

The repository's optimization checks compare matrix exponentials and block
limits with independent ODE solutions, stochastic moment formulas with exact
batch and short-trajectory enumeration, and stability thresholds with operator
eigenvalues on feasible binary supports. Rational path enumeration also checks
first-pass sampling equivalence and the effects of reusing observations.
Existing coordinate-descent and Gaussian integration checks validate the
explicit-penalty baselines. The pilot additionally checks conditional moments
and subgroup risks by exact path and support enumeration. These checks validate
the calculations, while the pilot provides a bounded comparison under its
declared tuning protocol. The practical study adds independent preprocessing,
transformed-coordinate update, held-out loss, selection, uncertainty, and
precision checks. Its categorical evidence remains specific to the task and
search grid. The novelty assessment must identify a conclusion beyond known
early-stopping, adaptive-gradient, and preconditioning theory before the paper
makes a practical recommendation.


\bibliographystyle{plainnat}
\bibliography{references}

\appendix
\section{Explicit-regularization baselines}
\label{sec:regularization}

The following results preserve the initial shrinkage investigation as baselines
for the optimization study. For fixed positive vectors $u,v$, consider the
original-scale estimator
\begin{equation}
 \label{eq:objective}
 \hat\beta = \arg\min_{\beta \in \R^p} \left\{ \frac12\beta^\top C\beta - z^\top\beta + t\sum_j u_j|\beta_j| + \frac r2\sum_j v_j\beta_j^2 \right\}, \qquad t,r
           \ge 0.
\end{equation}
Dividing centered feature $j$ by $s_j$ and transforming the fitted coefficient
back gives $u_j = s_j$ and $v_j = s_j^2$. This equivalence requires no
orthogonality. Allowing separate weights also covers penalty schemes such as
$u_j = v_j = C_{jj}$.

Weights and penalties are fixed with respect to $y$ throughout the distribution
calculations. Cross-validation and response-dependent penalty paths require
separate analysis. Write $\soft_k(x) = \sign(x)(|x| - k)_+$ for
soft-thresholding.

\begin{proposition}[Ridge in arbitrary dimension]
 \label{prop:ridge}
 Let $V = \diag(v)$ and $t = 0$. Then
 \begin{align*}
  \E\hat\beta           & = (C + rV)^{-1}C\beta^*,             \\
  \E\hat\beta - \beta^* & = -r(C + rV)^{-1}V\beta^*,           \\
  \Cov(\hat\beta)       & = \tau^2(C + rV)^{-1}C(C + rV)^{-1}.
 \end{align*}
 The mean and covariance formulas require only zero-mean errors with covariance
 $\sigma_\varepsilon^2 I$, rather than Gaussian errors.
\end{proposition}
\begin{proof}
 Differentiate the quadratic objective to obtain $\hat\beta = (C + rV)^{-1}z$,
 then take its mean and covariance.
\end{proof}
For $r > 0$ and $0 < k < 1$, the identity $(C + rV)^{-1}C = kI$ implies
$V = (1 - k)C/(kr)$. Thus a diagonal ridge penalty can give the same shrinkage
factor for every coefficient vector only when $C$ is diagonal.

\section{An exact two-feature representation}

Set
\[
 C = \begin{pmatrix}
      a & c \\
      c & b
     \end{pmatrix},\qquad M
   = C + r\diag(v)
   = \begin{pmatrix}
      d_1 & c   \\
      c   & d_2
     \end{pmatrix},\qquad \ell_j
   = tu_j,
\]
where $a,b > 0$ and $ab - c^2 > 0$. The KKT conditions are
\begin{equation}
 \label{eq:kkt}
 M\hat\beta + \ell \odot g = z,\qquad g_j
                           = \sign(\hat\beta_j)\ \text{if }\hat\beta_j
                           \ne 0, \quad |g_j|
                           \le 1\ \text{otherwise}.
\end{equation}
For $t > 0$, the solution has nine support/sign regions. On the empty support,
the conditions are $|z_j| \le \ell_j$ for both coordinates. On support $\{1\}$
with sign $h_1 \in \{-1,1\}$,
\[
 \hat\beta_1 = (z_1 - \ell_1h_1)/d_1,\qquad \hat\beta_2
             = 0, \quad h_1z_1
             > \ell_1,\quad |z_2 - c\hat\beta_1|
             \le \ell_2.
\]
Exchange the indices for support $\{2\}$. On the full support with signs
$h \in \{-1,1\}^2$,
\begin{equation}
 \label{eq:both-active}
 \hat\beta = M^{-1}(z - \ell \odot h),\qquad h_j\hat\beta_j > 0.
\end{equation}
These affine formulas and linear inequalities follow from \eqref{eq:kkt}. Strict
convexity makes them sufficient as well as necessary. Up to boundaries of
Gaussian probability zero, the regions partition score space.

\begin{proposition}[Exact selection integral]
 \label{prop:selection}
 Let $B_1(x) = \soft_{\ell_1}(x)/d_1$. Then
 \begin{equation}
  \label{eq:inactive}
  \hat\beta_2 = 0\quad \Longleftrightarrow \quad |z_2 - cB_1(z_1)| \le \ell_2.
 \end{equation}
 Write $\mu = C\beta^*$, $f_1$ for the $N(\mu_1,\tau^2a)$ density, and
 \[
  m(x) = \mu_2 + (c/a)(x - \mu_1),\quad q(x)
       = cB_1(x),\quad \omega
       = \tau\sqrt{b - c^2/a}.
 \]
 With $\Phi$ the standard normal distribution function,
 \begin{equation}
  \label{eq:selection}
  \Prob(\hat\beta_2 \ne 0) = \int_{\R} \left\{ \Phi\left( \frac{m(x) - q(x) - \ell_2}{\omega} \right) + \Phi\left( \frac{q(x) - m(x) - \ell_2}{\omega} \right) \right\}f_1(x)\,dx.
 \end{equation}
\end{proposition}
\begin{proof}
 The restricted minimizer $(B_1(z_1),0)$ solves the full problem precisely when
 the inactive KKT inequality holds. Conditional on $z_1 = x$, $z_2$ is normal
 with mean $m(x)$ and variance $\omega^2$. Add the two tail probabilities
 outside $[q(x) - \ell_2,q(x) + \ell_2]$ and integrate over $x$.
\end{proof}
Selection of feature 2 generally depends on the ridge penalty on feature 1,
through $d_1$, while its own ridge weight does not enter this event when other
parameters are held fixed. Numerical integration can split at $\pm\ell_1$.

The same regions yield exact integral representations for bias and covariance.
If the estimator is $A_kz + b_k$ on region $R_k$, define the unnormalized
Gaussian moments $p_k = \int_{R_k}f$, $m_k = \int_{R_k}zf$, and
$Q_k = \int_{R_k}zz^\top f$. Then
\begin{align}
 \E\hat\beta               & = \sum_k(A_km_k + b_kp_k),\label{eq:mean-regions}                                               \\
 \E\hat\beta\hat\beta^\top & = \sum_k\left( A_kQ_kA_k^\top + A_km_kb_k^\top + b_km_k^\top A_k^\top + p_kb_kb_k^\top \right).
\end{align}
Subtracting the appropriate mean products gives covariance. Truncation matters:
substituting $\E z$ into a single active-region formula does not give either the
unconditional mean or the mean conditional on that region.

\section{What binary structure preserves}
For binary features, write $q_j=\Prob(X_j=1)$ and $p_{11}=\Prob(X_1=X_2=1)$,
or use their empirical counterparts for a fixed design. Then
\[
 a=q_1(1-q_1),\quad b=q_2(1-q_2),\quad c=p_{11}-q_1q_2,
 \qquad \max(0,q_1+q_2-1)\le p_{11}\le\min(q_1,q_2).
\]
The conditional mean identities
\begin{equation}\label{eq:conditional-differences}
 \frac ca=\E(X_2\mid X_1=1)-\E(X_2\mid X_1=0),\qquad
 \frac cb=\E(X_1\mid X_2=1)-\E(X_1\mid X_2=0)
\end{equation}
give $|c|\le\min(a,b)$. Equality in either bound forces identical or
complementary columns, so both inequalities are strict when $C$ is positive
definite.

\begin{theorem}[Exact single-signal shrinkage]\label{thm:single}
In the noiseless two-binary-feature lasso, let $u=(a,b)$, $r=0$, and
$\beta^*=(B,0)$. For every $t\ge0$, the unique minimizer is
$\hat\beta=(\soft_t(B),0)$.
\end{theorem}
\begin{proof}
The first coordinate satisfies its scalar KKT condition. The inactive residual
has absolute value $|c(B-\soft_t(B))|\le |c|t\le bt$, so the second KKT
condition also holds. Strict convexity gives uniqueness.
\end{proof}
The result also holds with coordinates exchanged. Noise can still produce
false selections, so this is a statement about noiseless shrinkage.

For two active coefficients with signs $h$, \eqref{eq:both-active} gives
\begin{equation}\label{eq:binary-shrinkage}
 \hat\beta_1=\beta_1^*-\frac{t\{h_1-(c/a)h_2\}}{1-\rho^2},\qquad
 \hat\beta_2=\beta_2^*-\frac{t\{h_2-(c/b)h_1\}}{1-\rho^2},\qquad
 \rho=\frac{c}{\sqrt{ab}}.
\end{equation}
The resulting coefficients must have the stipulated signs. For a nested
design with $q_1=1/2$, $q_2=p_{11}=1/10$, we have
$(a,b,c)=(.25,.09,.05)$. At $t=.1$, the exact solutions are
\begin{center}
\begin{tabular}{lll}
\toprule
$\beta^*$ & Orthogonal, same marginals & Nested design \\
\midrule
$(1,1)$ & $(.9,.9)$ & $(.91,.95)$ \\
$(1,-1)$ & $(.9,-.9)$ & $(.865,-.825)$ \\
$(1,0)$ & $(.9,0)$ & $(.9,0)$ \\
\bottomrule
\end{tabular}
\end{center}
Thus correlation can increase or decrease shrinkage depending on the signs.
Both designs can be realized with integer cell counts, using sample sizes
that are multiples of twenty for the orthogonal comparison.

\begin{proposition}[No universal diagonal correction]\label{prop:no-diagonal}
If $c\ne0$ and $t>0$, no positive design-dependent weights $u$ make the
noiseless lasso equal $(\soft_t(\beta_1^*),\soft_t(\beta_2^*))$ for every
$\beta^*\in\R^2$.
\end{proposition}
\begin{proof}
For $(B,0)$ with $B>t$, the desired active KKT equation forces $u_1=a$.
Exchanging coordinates forces $u_2=b$. For $(B,B)$ with $B>t$, the desired
solution would then require $C(1,1)^\top=(a,b)^\top$, which forces $c=0$.
\end{proof}
If active signs were known, the weights $u_1=a+ch_1h_2$ and $u_2=b+ch_1h_2$
would instead produce $\hat\beta=\beta^*-th$ whenever $\beta^*$ has signs $h$
and $\min_j|\beta_j^*|>t$. Binary feasibility makes these weights positive.
Their dependence on unknown signs prevents a universal preprocessing rule.

\begin{proposition}[An approximation condition]\label{prop:bound}
Let $D=\diag(a,b)$ and $\eta=\max(|c|/a,|c|/b)<1$. On a valid noiseless
both-active sign region with variance weights,
\[
 \|\hat\beta-(\beta^*-th)\|_\infty\le t\frac{\eta}{1-\eta}.
\]
\end{proposition}
\begin{proof}
Write $C=D(I+F)$, where $\|F\|_\infty=\eta$. The Neumann series gives
$\|(I+F)^{-1}-I\|_\infty\le\eta/(1-\eta)$, and
$\hat\beta=\beta^*-tC^{-1}Dh$.
\end{proof}
Small conditional mean differences in \eqref{eq:conditional-differences}
therefore control the approximation. Small Pearson correlation alone does
not provide this control when variances differ greatly.

\section{A rare feature with vanishing correlation}
Keep $q_1\in(0,1)$ fixed and let $q_2\downarrow0$. Binary feasibility implies
$c=O(b)$, so $\rho\to0$, but $c/b$ need not vanish. Suppose $c/b\to\kappa$.
For nested binary features, $X_2=1$ implies $X_1=1$, and
\begin{equation}\label{eq:nested}
 c=q_2(1-q_1),\qquad \frac cb=\frac{1-q_1}{1-q_2}\longrightarrow1-q_1.
\end{equation}
With positive signals that remain active, \eqref{eq:binary-shrinkage} gives
\begin{equation}\label{eq:noiseless-limit}
 \hat\beta_1\longrightarrow\beta_1^*-t,\qquad
 \hat\beta_2\longrightarrow\beta_2^*-t(1-\kappa).
\end{equation}
For $q_1=1/2$, the rare coefficient loses $t/2$, compared with the orthogonal
loss $t$. Pearson correlation tends to zero in both comparisons.

This example is realizable along actual binary designs. Take $n_m=2m^2$,
with cell counts $(n_{11},n_{10},n_{01},n_{00})=(m,m^2-m,0,m^2)$ for
$m\ge2$. Then $q_1=1/2$, $q_2=1/(2m)$, and $n_mq_2=m\to\infty$.
Fix $t>0$ and $\beta_1^*,\beta_2^*>t$. Under independent Gaussian errors
with fixed variance, \eqref{eq:noiseless-limit} holds in probability for
variance-weighted lasso. Indeed, the positive-sign candidate equals the
noiseless candidate plus $C^{-1}(z-C\beta^*)$, whose covariance is
$\sigma_\varepsilon^2 C^{-1}/n_m$. Its largest coordinate variance is
$O((n_mq_2)^{-1})$, so its noise vanishes. The limiting coefficients are
strictly positive, and hence the candidate satisfies the active-sign
conditions with probability tending to one. This fixes the averaged penalty,
meaning $\lambda_1=n_mt$.

Ridge gives another explicit counterexample. With standard-deviation scaling,
$V=\diag(a,b)$, Proposition~\ref{prop:ridge} yields
\[
 \E\hat\beta_2=
 \frac{r(c/b)\beta_1^*+(1+r-\rho^2)\beta_2^*}{(1+r)^2-\rho^2}
 \longrightarrow\frac{\beta_2^*}{1+r}+
 \frac{r\kappa\beta_1^*}{(1+r)^2}.
\]
For $q_1=1/2$, $r=1$, and $\beta^*=(1,1)$, this mean tends to $5/8$,
whereas the orthogonal mean is $1/2$.

There is a distinct formal Gaussian experiment in which $\tau$ in
\eqref{eq:score-law} is fixed while the Gram matrix varies continuously.
Appendix~\ref{sec:formal} proves a rare-feature mean limit for this experiment.
It does not represent a fixed-size binary sample, whose empirical proportions
lie on a finite grid. Nor does it describe increasing samples with fixed
error variance, since those have $\tau_n\to0$. These distinctions matter
because thresholding, rare counts, and penalties can interact differently
under different limiting regimes.

\section{A formal Gaussian Gram-matrix limit}\label{sec:formal}
\begin{theorem}\label{thm:formal}
Hold $a,t,\tau>0$, $r\ge0$, and $\beta^*\in\R^2$ fixed. Let $b\downarrow0$
and $c/b\to\kappa$, and use $u=(a^\delta,b^\delta)$ and
$v=(a^{2\delta},b^{2\delta})$ for $\delta\ge1/2$. Define
\[
 B_1=\frac{\soft_{ta^\delta}(a\beta_1^*+\xi_1)}{a+ra^{2\delta}},
 \quad \xi_1\sim N(0,\tau^2a),\qquad
 L=\beta_2^*+\kappa(\beta_1^*-\E B_1).
\]
Under \eqref{eq:score-law}, the limits are
\begin{align*}
 \lim\E\hat\beta_2&=
 \begin{cases}2\Phi(-t/\tau)L/(1+r),&\delta=1/2,\\L,&\delta>1/2,\end{cases}\\
 \lim\Prob(\hat\beta_2\ne0)&=
 \begin{cases}2\Phi(-t/\tau),&\delta=1/2,\\1,&\delta>1/2,\end{cases}\\
 \lim b\Var(\hat\beta_2)&=
 \begin{cases}\Var(\soft_t(\tau Z))/(1+r)^2,&\delta=1/2,\\
 \tau^2,&\delta>1/2,\end{cases}\qquad Z\sim N(0,1).
\end{align*}
\end{theorem}
\begin{proof}
Couple the score errors by $\xi_2=(c/a)\xi_1+
\tau\sqrt{b-c^2/a}\,Z$, with $Z$ independent of $\xi_1$.
The coordinate KKT equations and the Lipschitz property of soft-thresholding
give $\|\hat\beta_2\|_{L^2}=O(b^{-1/2})$ and
$\|\hat\beta_1-B_1\|_{L^2}=O(\sqrt b)$. For the first bound, combine
$d_j\|\hat\beta_j\|_{L^2}\le\|z_j\|_{L^2}+|c|\|\hat\beta_{3-j}\|_{L^2}$
and use $d_2-c^2/d_1\asymp b$; for the second, apply the Lipschitz bound to
the first coordinate equation. Replacing $\hat\beta_1$ by $B_1$ in the second
coordinate changes its estimate by $O_{L^2}(\sqrt b)$. The replacement is
\[
 \widetilde B_2=
 \frac{\soft_{tb^\delta}(\tau\sqrt{b-c^2/a}\,Z+bU_b)}{b+rb^{2\delta}},
 \qquad U_b=\beta_2^*+(c/b)(\beta_1^*-B_1+\xi_1/a).
\]
Here $U_b$ is independent of $Z$, converges in $L^2$, and has mean tending
to $L$. Define
\[
 F_b(h)=\E_Z\soft_{tb^\delta}(\tau\sqrt{b-c^2/a}\,Z+h).
\]
Symmetry gives $F_b(0)=0$ and differentiation gives
\[
 F_b'(h)=\Prob_Z(|\tau\sqrt{b-c^2/a}\,Z+h|>tb^\delta),\qquad
 \frac{F_b(bU_b)}b=U_b\int_0^1F_b'(sbU_b)\,ds.
\]
The integrand tends to $2\Phi(-t/\tau)$ when $\delta=1/2$ and to one
otherwise. It is bounded by one, and the $L^2$ bound on $U_b$ gives uniform
integrability. Taking expectations and dividing by $1+rb^{2\delta-1}$
proves the mean limit.

The same representation gives $\sqrt b\hat\beta_2\to
\soft_t(\tau Z)/(1+r)$ in $L^2$ when $\delta=1/2$, and $\sqrt b\hat\beta_2\to
\tau Z$ in $L^2$ otherwise. Both limiting variables have mean zero, proving
the variance limits. For selection, use \eqref{eq:inactive}; its residual
divided by $\sqrt b$ converges in distribution to $\tau Z$, while its
threshold tends to $t$ or zero. The limiting boundary has probability zero.
\end{proof}
This formal experiment keeps score noise $\tau$ fixed. A statistical theorem
for actual binary samples must replace it with a specified joint limit for
$n$, prevalence, error variance, and penalties.

\end{document}
